\subsection{Um Pouco de �lgebra Linear}

Fa�a a seguinte brincadeira. Na equa��o de onda, o operador diferencial (fazendo c = 1) � do seguinte tipo
%
%
%
\[D = \deli{^2}{x^2} - \deli{^2}{t^2}\]
%
%
%
Vamos criar um polin�mio a partir desse operador, use as substitui��es:
%
%
%
\[\deli{}{x} \to ix \qquad \qquad \deli{}{t} \to it \]
%
%
%
Ficamos com:
%
%
%
\[p(x,t) = t^{2} - x^{2} \]
%
%
%
Tenho certeza de que voc� j� percebeu. Isso mesmo! � a m�trica de Minkowski! O Eintein, ao exigir que a equa��o de onda fosse preservada (invariante), indiretamente, estava calculando o grupo de Lorentz! Voltando para o polin�mio $p(x,t)$, temos que as curvas de n�vel para $p(x,y) > 0$ s�o hip�rboles, ent�o o que voc� deduziu no final, podemos chamar de rota��o hiperb�lica.
